\documentclass[landscape]{slides}
\usepackage[margin=1cm,nohead]{geometry}
\usepackage{amsmath,amssymb}

\usepackage{xcolor}
\usepackage{graphicx}
\usepackage{subcaption}
\usepackage{enumitem}

\usepackage{array} 

\def\myref#1{\textcolor{magenta}{\small #1}}
\def\myemph#1{\textcolor{blue}{\em #1}}
\def\bigemph#1{\textcolor{red}{\em #1}}
\def\formula#1{{\color{black}$#1$}}

\def\clsp#1{\vspace{-#1}}


\def\Rnum{\mathbb{R}}
\def\Cnum{\mathbb{C}}
\def\txtint{{\textstyle\int}}
\def\bigint{{\displaystyle\int}}
\def\sgn{\text{sgn}}
\def\const{\text{const.}}

\def\smallbinom#1#2{{\textstyle \binom{#1}{#2}}}
\def\smallsum{\textstyle\sum}
\def\T{{\rm T}}
\def\i{\hat{i}}

\def\scal{{\rm scal.}}
\def\trans{{\rm trans.}}
\def\triv{{\rm triv.}}
\def\Jsp{\mathrm{J}^{(\infty)}}
\def\Esp{\mathcal{E}}

\def\Lagr{{\mathcal L}}

\def\longit{{(u)}}
\def\transv{{(v)}}



\definecolor{pakistangreen}{rgb}{0.0, 0.4, 0.0}
\definecolor{officegreen}{rgb}{0.0, 0.5, 0.0}
\definecolor{lasallegreen}{rgb}{0.03, 0.47, 0.19}



\def\X#1{{\mathrm X}_{#1}}
\def\Y#1{{\mathrm Y}_{#1}}
\def\hatX#1{\hat{\mathrm X}_{#1}}
\def\hatY#1{\hat{\mathrm Y}_{#1}}
\def\pr{\mathrm {pr}}

\def\Ref#1{Ref.\cite{#1}}



\tolerance=99999


\title{Nonlinear wave equation for an elastic string: \\ derivation, symmetries, conserved quantities}
\author{Mathew Alex 7652597
\\ 
Supervisor: Prof. Stephen Anco
\\
M.Sc. Physics Thesis Presentation
}
\date{July 31, 2025}

\begin{document}
\maketitle

\parskip=3pt


\newpage
%intro-1

\begin{center}
{\bf  Linear Waves}
\end{center}

Linear waves in 1D elastic media have several important physical\\ applications:
\begin{enumerate}
\item 
vibrations of an elastic string e.g.\ violin or piano
\item 
sound waves e.g.\ pipe organ or flute
\item 
seismic waves e.g.\ earth tremor
\end{enumerate}
\hfil

$v(x,t) = \text{transverse displacement}$

equilibrium is $v(x,t)\equiv 0$

\includegraphics[width=0.8\textwidth]{Presentation-diagram1.pdf}

Governing equation for linear waves
\[
v_{tt} - c^2 v_{xx} = 0, 
\quad
c=\const
\]

\newpage
%intro-2

Physical conditions under which the linear wave equation holds for the motion of an elastic string:
\begin{enumerate}
\item 
constant tension 
\item 
small angle approximation ($ |v_x| \ll 1 $)
\item 
no longitudinal motion
\end{enumerate}
\hfil

For motion that involves \emph{large deformations} from equilibrium, \\
\myemph{all three conditions no longer hold}. 
\\
    
Waves in any real elastic material will exhibit some \emph{longitudinal motion}
(i.e.\ compression/rarefaction will occur). 
\\

\begin{center}
\myemph{We want to study correct physical wave equation \\ for an elastic string without restrictions.}
\end{center}

\newpage
%intro-3

\begin{center}
\textbf{Nonlinear wave equation for  1D elastic string}
\end{center}
\hfil

\begin{itemize}
\item 
dynamical state of an elastic string can involve motion 
in both transverse directions and longitudinal direction 
\item 
physical elastic string will have finite length with boundary conditions at the ends
\item 
in a very long elastic string, motion in interior does not depend on boundary conditions 
\end{itemize}
\hfil

For mathematical simplicity:\\
consider only free motion without boundary conditions
\\

$v,w=\text{ transverse displacements}$; $u=\text{ longitudinal displacement}$ 

equilibrium is $v,w\equiv 0$ and $u\equiv 0$
\\

\includegraphics[width=0.8\textwidth]{Presentation-diagram2.pdf}

\newpage

\includegraphics[width=0.8\textwidth]{Presentation-diagram3.pdf}

Elastic string wave equations

\begin{align*}
u_{tt} & = c^2 u_{xx} + ({\tilde c}^2 -c^2)\big( (1+u_x)\big/\sqrt{(1+u_x)^2 + v_x^2 + w_x^2} \big){}_x
\\
v_{tt} & = c^2 v_{xx} + ({\tilde c}^2 -c^2)\big( v_x\big/\sqrt{(1+u_x)^2 + v_x^2 + w_x^2} \big){}_x
\\
w_{tt} & = c^2 w_{xx} + ({\tilde c}^2 -c^2)\big( w_x\big/\sqrt{(1+u_x)^2 + v_x^2 + w_x^2} \big){}_x
\end{align*}
where $c,\tilde c=\const > 0$

\newpage
%survey-1
\begin{center}
\textbf{Survey of literature on derivations}
\end{center}

\textbf{First derivation:} W.F. Osgood \myref{1925, ``Advanced Calculus'' (book)}

\textbf{Widely-cited paper:} G.F. Carrier \myref{1945, Q. J. Appl. Math.}

\begin{itemize}
\item
applied Newton's law framework
\item
restricted motion to a plane 
\item
assumes linear stress-strain relation (Hooke's law) 
\end{itemize}
\hfil

Longitudinal dynamical variable $u$ and transverse dynamical variable $v$\\
$\Longleftrightarrow$ $w=0$ in the nonlinear wave equations 
\\

Stress $\sigma = E\lambda$, 
where $\lambda = \sqrt{(1+u_x)^2 + v_x^2} -1$ is strain (net stretch)\\
and $E$ is elastic modulus (constant)
\\

Dynamical tension $T = T_0 + A\sigma$,
where $T_0$ is tension in equilibrium\\
and $A$ is cross-section area (constant)

\[\begin{aligned}
\rho A u_{tt} & = \big( (1+u_x)T/(\lambda +1) \big){}_x, 
\quad
\rho A v_{tt} & = \big( v_xT/(\lambda +1) \big){}_x
\end{aligned}\]
where $\rho$ is density (constant)

\newpage
%survey-2
\textbf{Subsequent major paper:} R. Narasimha \myref{1968, J. Sound and Vibration}

Similarly to Osgood/Carrier
\begin{itemize}
\item
applied Newton's law framework 
\item
assumes linear stress-strain relation (Hooke's law) 
\end{itemize}
\hfil

Differences
\begin{itemize}
\item
used Lagrangian variables (position and time for motion of particles comprising the elastic string)
\item
\myemph{no restriction} on transverse motion 
\item
emphasizes that, in general, \myemph{longitudinal mode is excited} ($u\not\equiv 0$)
\item
points out that \myemph{transverse motion in a plane may be unstable}
\end{itemize}
\hfil

When the equations are re-written in the ``lab'' variables $x$ and $t$, 
they are equivalent to the equations for $(u,v,w)$

\newpage
%survey-3

\textbf{Nonlinear elasticity approach:} S. Antman \myref{2005, ``Noninear Problems of Elasticity" (book)}

\begin{itemize}
\item
uses conservation of momentum and angular momentum \\
in vector variables
\item
no restriction on transverse motion
\item 
considers general \myemph{nonlinear stress-strain} relation
\end{itemize}
\hfil

position vector of material points $\vec{r}(t,x) = (x+u,v,w)$
\\

elastic force vector $\vec{F}^e$
\\

Conservation of momentum states
\[
\rho A \vec{r}_{tt} = \vec{F}^e_x
\]

Conservation of angular momentum leads to the condition 
$\vec{F}^e \times \vec{r}_x=0$ \\
(called ``perfect flexibility'') \\
$\Longrightarrow$
\[
\vec{F}^e = K(\lambda) \widehat{r_x}
\]
where $\widehat{r_x} = |\vec{r_x}|^{-1} \vec{r_x}$ is unit tangent vector,
and $K(\lambda)$ is called the stress-strain function.

\newpage

Stress-strain function $K(\lambda)$ must satisfy two physical conditions
\begin{itemize}
\item
$K$ increases with $\lambda$: string lengthens only with increasing $K$
\item
$K\to\infty$ as $\lambda\to\infty$: an infinitely large stretch occurs only by an infinite tension
\end{itemize}
\hfil

Equation of motion for $\vec{r}(t,x)$ has explicit component form
\begin{align*}
u_{tt} & =  \big( k(\lambda) (1+u_x) \big){}_x
\\
v_{tt} & = \big( k(\lambda) v_x \big){}_x
\\
w_{tt} & = \big( k(\lambda) w_x \big){}_x
\end{align*}
with
\[
k(\lambda) = \frac{K(\lambda)}{(\lambda+1)\rho A},
\quad
\lambda = \sqrt{(1+u_x)^2 + v_x^2 + w_x^2} -1
\]
\hfil

\begin{itemize}
\item
Linear function $K(\lambda) = T_0 + AE\lambda$
yields Narasimha's equations for elastic string
\item
Nonlinear functions $K(\lambda)$ describe nonlinearly elastic string\\
(i.e.\ nonlinear tension force)
\end{itemize}

\newpage
%survey-4

\textbf{Continuum mechanics framework:}  
Serrano et al. \myref{2019, Internat. J. Solids and Structures}

\begin{itemize}
\item
uses \myemph{action principle} in terms of displacement variables
\item
allows for general transverse motion
\item 
based on stored potential energy which is a function of strain tensor
\item
equations of motion given by variational derivative of action principle
\end{itemize}
\hfil

equilibrium is $\vec{X} = (x,0,0)$ (called ``reference configuration'')
\\

$\vec{r} = (x+u,v,w)$ position of material elements comprising string\\
(called ``current configuration'')
\\

displacement is $\vec{R} = \vec{r} - \vec{X} = (u,v,w)$
\\

kinetic energy of material element $\tfrac{1}{2}\rho A |\vec{R}_t|^2 dx$
\\

a material element $d\vec{X}$ is stretched to  
$d\vec{r} 
= 
\mathbf{F} \cdot d\vec{X}$,
where 
$\mathbf{F} = \partial_{\vec{X}} \vec{r}$ is the deformation gradient tensor
\\

stored potential energy of a material element is
$W(\text{strain tensor})dx$,
where the strain tensor, quantifiying the deformation, vanishes when $\mathbf{F} \in \text{SO}(3)$
\\

some obvious choices for the strain tensor are 
$(\mathbf{F}^\top \mathbf{F})^p - \mathbf{I} 
= (\mathbf{F}^\top \mathbf{F}) \cdots (\mathbf{F}^\top \mathbf{F}) - \mathbf{I}, p \neq 0$,
and $\ln \mathbf{F}^\top \mathbf{F}$
\\

physically, $\mathbf{F}^\top \mathbf{F}$ gives the norm of the stretched material element, 
$d\vec{r} \cdot d\vec{r} = (d\vec{X})^\top \mathbf{F}^\top \mathbf{F} d\vec{X}$
\\

for strings,
$\mathbf{F} = |\vec{r}_x| (\widehat{r_x} \otimes \hat{i})$
and
$\mathbf{F}^\top \mathbf{F} = |\vec{r}_x|^2 (\hat{i} \otimes \hat{i}) 
= (1 + \lambda)^2 (\hat{i} \otimes \hat{i})$,
where 
$\lambda = |\vec{r}_x|-1
= \sqrt{(1 + u_x)^2 + v_x^2 + w_x^2} - 1$ 
is the strain (net stretch)
\\

effectively, the stored potential energy can be defined to be 
$W(\lambda) dx$

\newpage

Action principle is total kinetic energy minus total potential energy 
\[
S = \int_{-\infty}^{\infty} \int_{-\infty}^{\infty} \big( \tfrac{1}{2}\rho A |\vec{R}_t|^2 -W(\lambda) \big) dx dt
\]

$\Longrightarrow$ variational derivative yields (material) equations of motion
\[
\frac{\delta S}{\delta \vec{R}} = -\rho A \vec{R}_{tt} + \big( (W'(\lambda)/(\lambda+1)) \vec{r}_x \big){}_x =0
\]

$\Longrightarrow$ nonlinear elastic string equations 
\[
\vec{R}_{tt}  = ( k(\lambda) \vec{r}_x )_x ,
\quad
k(\lambda) = \frac{W'(\lambda)}{(\lambda+1)\rho A},
\quad
\lambda = |\vec{r}_x|-1,
\quad
\vec{r} = x\,\hat{i} +\vec{R}
\]
where $W'(\lambda) = K(\lambda)$ is the stress-strain function. 
\\

\begin{itemize}
\item
Potential energy is typically a quadratic function of strain
\[
W(\lambda) = \tfrac{1}{2} c^2 (\lambda+1)^2 + ({\tilde c}^2 - c^2) \lambda
\]
$\Longrightarrow$ linear stress-strain $K(\lambda) = c^2 \lambda + {\tilde c}^2$
\item
Reduces to the Carrier/Narasimha equations for $\vec{R}=(u,v,w)$ 
\end{itemize}

\newpage
% variational structure
\begin{center}
\textbf{ Lagrangian and Hamiltonian of elastic string equations}
\end{center}
\begin{itemize}
\item
Lagrangian 
\[
L[\vec{R}] = \tfrac{1}{2}\rho A |\vec{R}_t|^2 - W(\lambda)
\]
\[
\frac{\delta L}{\delta \vec{R}} = -\rho A \vec{R}_{tt} + \big( (\underbrace{W'(\lambda) / (\lambda + 1)}_{\displaystyle \rho A k(\lambda) }) \vec{r}_x  \big){}_x = 0
\]
\item
Hamiltonian, $H$, can be obtained from $L$ via a Legendre transform using the momentum 
\[
\vec{p} = \dfrac{\partial L}{\partial \vec{R_t}} = \rho A \vec{R}_t
\]
$\Longrightarrow$
\[
H[\vec{p},\vec{R}] = \vec{p} \cdot \vec{R_t} - L 
= \frac{|\vec{p}|^2}{2 \rho A} + W(\lambda)
\]
Hamilton's equations of motion:
\[
\vec{R}_t = \frac{\delta H}{\delta \vec{p}} = \frac{1}{\rho A}\vec{p}
\quad\text{  and }\quad
\vec{p}_t = -\frac{\delta H}{\delta \vec{R}} 
= \rho A \big( k(\lambda) \vec{r}_x  \big){}_x 
\]
\item
These equations are evolution PDEs which have a unique solution\\
for given initial data $(u,v,w)\big|_{t=0}$ and $(u_t,v_t,w_t)\big|_{t=0}$. 
\end{itemize}


\newpage
%symms-1
\begin{center}
\textbf{Point symmetries}
\end{center}

\begin{itemize}
\item 
Transformations acting in the space of kinematic variables $(x,t,u,v,w)$
under which the form of the wave equations is preserved 
\item 
The wave equations manifest some evident point symmetries: \\
\emph{time and space translations}; \emph{rotations in the transverse $(v,w)$ plane}
\end{itemize}
\hfil

\begin{center}
\myemph{Goal: Find all non-manifest point symmetries}
\end{center}
\hfil

Consider one-parameter group of transformations 
\begin{align*}
x \to \tilde{x} & = x + \epsilon \xi(x,t,u,v,w) + O(\epsilon^2) \\
t \to \tilde{t} & = t + \epsilon \tau(x,t,u,v,w) + O(\epsilon^2) \\
u \to \tilde{u} & =  u + \epsilon \eta^u(x,t,u,v,w) + O(\epsilon^2) \\
v \to \tilde{v} & = v + \epsilon \eta^v(x,t,u,v,w) + O(\epsilon^2) \\
w \to \tilde{w} & = w + \epsilon \eta^w(x,t,u,v,w) + O(\epsilon^2) 
\end{align*}
where $\epsilon = 0$ gives the identity transformation. 

\begin{itemize}
\item
Need to know how transformations act on functions of $(t,x)$, 
because the string equations contain derivatives of $(u,v,w)$
\item
Simplest way is by working with infinitesimal transformations 
\[\begin{aligned}
(u(t,x),v(t,x),w(t,x)) \to & (\tilde{u}(t,x),\tilde{v}(t,x),\tilde{w}(t,x)) \\
& = (u(t,x),v(t,x),w(t,x)) \\&\qquad\quad
+ \epsilon\, (P^u(t,x),P^v(t,x),P^w(t,x)) + O(\epsilon^2)
\end{aligned}\]
which is given by 
\[
P^u = \eta^u - \xi u_x - \tau u_t \equiv \delta u,
\
P^v = \eta^v - \xi v_x - \tau v_t \equiv \delta v,
\
P^w = \eta^w - \xi w_x - \tau w_t \equiv \delta w
\]
\myref{Bluman \&  Anco 2002, ``Symmetry and Integration Methods for Differential Equations" (book) }
\item
$\delta$ denotes infinitesimal change given by the $\epsilon$-derivative at $\epsilon=0$
\end{itemize}
\hfil

\underline{Example}: space translation $x\to \tilde{x} = x +\epsilon$ 
has $\xi =1$, $\tau=0$, $\eta=0$

Transformation on variables gives
\[\begin{aligned}
& \tilde{u} = u 
&&\text{ since no change in $u$ }\\
& =\tilde{u}(\tilde{x},\tilde{t}) = u(x,t) =u(\tilde{x}-\epsilon,\tilde{t}) 
&&\text{ by $x=\tilde{x} -\epsilon$ and $t=\tilde{t}$ }
\end{aligned}\]
$\Longrightarrow$
$u(x,t) \to \tilde{u}(x,t) = u(x-\epsilon,t)$
equivalent transformation on functions
\\

Infinitesimal change in $u(t,x)$ is given by $P^u = \dfrac{\partial u(x-\epsilon,t)}{\partial \epsilon}\big|_{\epsilon=0} = -u_x$

\newpage
%symms-2

\textbf{How to find infinitesimal point symmetries}

Preservation of the form of the elastic string equations 
\[
\vec{R}_{tt} = ( k(\lambda) (\hat{i} + \vec{R}_x) ){}_x
\]
under a symmetry means that the symmetry transformation maps\\
solutions into solutions
\\

$\Longrightarrow$
Infinitesimal change in the form of the equations must vanish\\
for all solutions:
\[
\delta \left( \vec{R}_{tt} - ( k(\lambda) (\hat{i} + \vec{R}_x) ){}_x \right) 
=0
\quad\text{ when }\quad
\vec{R}_{tt} = ( k(\lambda) (\hat{i} + \vec{R}_x) ){}_x
\]
where $\delta \vec{R} = \vec{P} = (P^u,P^v,P^w)$.
\\

Substitute
$\delta \vec{R}_{tt} =\partial_t^2 \vec{P}$
and
$\delta \vec{R}_{x} =\partial_x \vec{P}$
and
$\delta k(\lambda) = \dfrac{k'(\lambda)}{\lambda +1} (\hat{i}+\vec{R}_x)\cdot \partial_x \vec{P}$
\\
where $\partial$ denotes the derivative acting by chain rule
\\
$\Longrightarrow$
\begin{align*}
& \partial_t^2 \vec{P} 
- \partial_x \left( 
\frac{k'(\lambda)}{\lambda + 1} \big( (\hat{i}+\vec{R}_x)\cdot \partial_x \vec{P} \big)(\hat{i}+\vec{R}_x)
+ k(\lambda) \partial_x \vec{P} \right)
=0
\\
&\quad\text{ when }\quad
\vec{R}_{tt} = ( k(\lambda) (\hat{i} + \vec{R}_x) ){}_x
\end{align*}
called \myemph{symmetry determining equations}

\newpage
%symms-3

More compact tensorial expression for the determining equations 
\[
\partial_t^2 \vec{P} \big|_{\vec{R}_{tt} = ( k(\lambda) (\hat{i} + \vec{R}_x) ){}_x}
= \partial_x \left( \mathbf{K} \cdot \partial_x \vec{P} \right)
\]
where
\[
\mathbf{K} = 
\frac{k'(\lambda)}{\lambda + 1} 
(\hat{i}+\vec{R}_x)\odot (\hat{i}+\vec{R}_x)
+ k(\lambda) \mathbf{I} 
\]
\hfil

\begin{itemize}
\item
For point symmetries
\[
\vec{P} = \vec{\eta} - \xi \vec{R}_x -\tau \vec{R}_t
\]
where $\vec{\eta}=(\eta^u,\eta^v,\eta^w)$, $\xi$, $\tau$ 
are unknown functions of $t,x,\vec{R}=(u,v,w)$
\item
By inspection:\\
\myemph{time translation} $\tau =1$, $\xi=0$, $\vec{\eta}=0$: \\
$\tilde{t}=t+\epsilon$

\myemph{space translation} $\xi =1$, $\tau=0$, $\vec{\eta}=0$: \\
$\tilde{x}=x+\epsilon$

\myemph{transverse rotation} $\vec{\eta} =(0,w, -v)$, $\xi =0$, $\tau=0$: \\
$\begin{pmatrix} \tilde{v} \\ \tilde{w} \end{pmatrix}
= \begin{pmatrix} 
\cos(\epsilon) & \sin(\epsilon) \\
-\sin(\epsilon) & \cos(\epsilon) 
\end{pmatrix}
\begin{pmatrix} v \\ w \end{pmatrix}$\\

\myemph{constant shifts} in $(u,v,w)$
\end{itemize}



\newpage
%symms-4

\textbf{Maple computation to find the point symmetries}
\hfil\\

Main steps:
\begin{enumerate}
\item 
Set up the general form for the elastic string equations 
and the general form for a point symmetry.
\\
\item
Compute the determining equations for point symmetries.
\\
\item
Bring the determining equations to simplest form such that 
linear combinations and cross-differentiations do not yield 
any equations of lower order in derivatives.\\

Different cases may arise for the resulting simplified equations. 
\\
\item
Solve in each case the system of equations for the unknowns.
\\
\end{enumerate}
\hfil


(Show worksheet)

\newpage
%symms-5
\begin{center}
\textbf{List of solutions modulo translations and transverse rotations}
\end{center}

\begin{itemize}
\item
Maple computation yields a classification of all point symmetries 
\item
Comprises three separate cases depending on the form of $k(\lambda)$
\end{itemize}
\hfil

\begin{enumerate}
\item 
Arbitrary $k(\lambda)$

6 point symmetries 

\begin{enumerate}[label=\roman*.]
\item 
\myemph{dilation}:  
$ (\tilde x,\tilde t,\tilde u,\tilde v,\tilde w) = e^{\epsilon} (x,t,u,v,w)$
\item
\myemph{time-dependent shift in $u$}:  
$\tilde u =u + \epsilon t$
\item
\myemph{time-dependent shift in $v$}:  
$\tilde v =v + \epsilon t$
\item
\myemph{time-dependent shift in $w$}:  
$\tilde w =w + \epsilon t$
\item 
\myemph{rotations in the $(x+u,v)$ plane}: \\
$\begin{pmatrix} \tilde x + \tilde u \\ \tilde v  \end{pmatrix}
= 
\begin{pmatrix}
\cos(\epsilon) & \sin(\epsilon) \\
-\sin(\epsilon) & \cos(\epsilon)
\end{pmatrix}
\begin{pmatrix} x + u \\ v \end{pmatrix}$
\item 
\myemph{rotations in the $(x+u,w)$ plane}: \\
$\begin{pmatrix} \tilde x + \tilde u \\ \tilde w  \end{pmatrix}
= 
\begin{pmatrix}
\cos(\epsilon) & \sin(\epsilon) \\
-\sin(\epsilon) & \cos(\epsilon)
\end{pmatrix}
\begin{pmatrix} x + u \\ w \end{pmatrix}$
\end{enumerate}

\newpage

\item 
$k(\lambda) = k_0 (1+\lambda)^{p}$,
\quad
$p \neq -4, -1, 0$, 
\quad
$k_0 = \const \neq 0$

1 extra point symmetry

\begin{enumerate}[label=\roman*.]
\item 
\myemph{scaling}: 
$\tilde t = e^{\epsilon} t$,
$\tilde u = e^{-\frac{2}{p}\epsilon} (x+u) - x$,
$\tilde v = e^{-\frac{2}{p} \epsilon} v$,
$\tilde w = e^{-\frac{2}{p} \epsilon} w$\\

invariants are $x$, $t^{\frac{2}{p}} v$, $t^{\frac{2}{p}} w$, $t^{\frac{2}{p}} (x+u)$ 
\end{enumerate}
\hfil

\item
$k(\lambda) = k_0 (1+\lambda)^{-4}$,
\quad
$k_0 =\const\neq 0$

%2 extra point symmetries
1 extra point symmetry

\begin{enumerate}[label=\roman*.]
% \item 
% \myemph{scaling}: 
% $\tilde t = e^{\epsilon} t$,
% $\tilde u = e^{\frac{1}{2}\epsilon} (x+u) - x$,
% $\tilde v = e^{\frac{1}{2} \epsilon} v$,
% $\tilde w = e^{\frac{1}{2} \epsilon} w$\\

% invariants are $x$, $\dfrac{v}{\sqrt{t}}$, $\dfrac{w}{\sqrt{t}}$, $\dfrac{x+u}{\sqrt{t}}$, 
%$t^{-\frac{1}{2}} v$, $t^{-\frac{1}{2}} w$, $t^{-\frac{1}{2}} (u+x)$ 
\item 
\myemph{time-dependent (conformal) scaling}:\\ 
$\tilde t = \dfrac{t}{1 - \epsilon t}$,
$\tilde u = \dfrac{u + \epsilon xt}{1 - \epsilon t}$, 
$\tilde v =\dfrac{v}{1 - \epsilon t}$,
$\tilde w = \dfrac{w}{1 - \epsilon t}$
\\

invariants are $x$, $\dfrac{v}{t}$, $\dfrac{w}{t}$, $\dfrac{x + u}{t}$
\end{enumerate}

\end{enumerate}

\newpage
\begin{center}
\textbf{Physical meaning of the $k(\lambda)$ in each case}
\end{center}
\hfil\\
The stress-strain function, $K(\lambda)$, in each case can be found using  $k(\lambda) = \frac{K(\lambda)}{(\lambda+1)\rho A}$: 
\\
\begin{enumerate}
\item 
Arbitrary $k(\lambda) \implies  \text{ arbitrary } K(\lambda)$

\item 
$k(\lambda) = k_0 (1+\lambda)^{p}$,
\quad
$p \neq -4, -1, 0$, 
\quad
$k_0 = \const \neq 0$ \\
$\implies K(\lambda) = K_0 (1+\lambda)^{p+1}, K_0 = \const \neq 0$

\item
$k(\lambda) = k_0 (1+\lambda)^{-4}$,
\quad
$k_0 =\const\neq 0$ \\
$\implies K(\lambda) = K_0 (1+\lambda)^{-3}, K_0 = \const \neq 0$

\end{enumerate}
\hfil 

\newpage

\begin{center}
\textbf{Conservation laws}
\end{center}

\begin{itemize}
\item 
By applying Noether's theorem 
\myref{Bluman, Cheviakov, \&  Anco 2010, ``Applications of symmetry methods to partial differential equations" (book) },
we can use the infinitesimal symmetries $\delta \vec{R} = \vec{P}$ 
to find the conservation laws admitted by the elastic string equations. 
\item
\myemph{Variational symmetries} are transformations that preserve the action.
Equivalently, they change the Lagrangian $L$ by a total space-time derivative, 
\[\label{variational_symmetry_condition1}
\delta_{\vec{P}} L \equiv \partial_t A + \partial_x B
\]
where $A$ and $B$ are functions of the dynamical variables and their derivatives.
\item
The \myemph{variational derivative annihilates total space-time derivatives}. 
Using this property, 
we can find the variational symmetries 
by checking the condition 
\[
\frac{\delta (\delta_{\vec{P}} L) }{\delta \vec{R}} = 0
\]
which holds off of solutions. 
\newpage
\item
For any variation $\delta \vec{R} = \vec{P}$, the infinitesimal change in the Lagrangian is 
\[\label{L_variation} \begin{aligned}
\delta_{\vec{P}} L 
&= 
\partial_t \vec{P} \cdot \partial_{\vec{R}_t} L
+ \partial_x \vec{P} \cdot \partial_{\vec{R}_x} L 
\\
&= 
\partial_t (\vec{P} \cdot \partial_{\vec{R}_t} L)
- \vec{P} \cdot \partial_t \partial_{\vec{R}_t} L
+ \partial_x (\vec{P} \cdot \partial_{\vec{R}_x} L)
- \vec{P} \cdot \partial_x \partial_{\vec{R}_x} L 
\\
&= 
\partial_t (\vec{P} \cdot \partial_{\vec{R}_t} L)
+ \partial_x (\vec{P} \cdot \partial_{\vec{R}_x} L)
- \vec{P} \cdot \frac{\delta L}{\delta \vec{R}}  
\end{aligned}\]
Rearranging this expression, we obtain
\[\label{variational_symmetry_condition2}\begin{aligned}
\vec{P} \cdot \frac{\delta L}{\delta \vec{R}}
&=
\partial_t (\vec{P} \cdot \partial_{\vec{R}_t} L)
+ \partial_x (\vec{P} \cdot \partial_{\vec{R}_x} L)
- \delta_{\vec{P}} L 
\\
&=
\partial_t (\vec{P} \cdot \partial_{\vec{R}_t} L - A)
+ \partial_x (\vec{P} \cdot \partial_{\vec{R}_x} L - B)
\end{aligned}\]
for variational symmetries.
\item
When equation evaluated on solutions, $\frac{\delta L}{\delta \vec{R}} = 0$,  
the space-time total derivative on the right-hand side evaluates to zero, yielding a conservation law:
\[
\vec{P} \cdot \frac{\delta L}{\delta \vec{R}} = \partial_t T + \partial_x X = 0
\] 

\newpage

where 
\[
T = \vec{P} \cdot \partial_{\vec{R}_t} L - A
\]
is the \myemph{conserved density}, 
and 
\[
X = \vec{P} \cdot \partial_{\vec{R}_x} L - B
\]
is its \myemph{flux}. 
\item
To find $T$ and $X$ explicitly, 
we repeatedly integrate by parts in $ \vec{P} \cdot \frac{\delta L}{\delta \vec{R}}$ 
until the expression takes the form of a space-time total derivative.
\end{itemize}
\hfil

\newpage

\textbf{Maple computation to find the conservation laws}


The main steps for computing conservation laws in Maple are:
\begin{enumerate}
\item
In each case, determine the conditions on the symmetry generator such that the variational derivative annihilates the dot product of the symmetry generator and the string equations.
\item
Substitute these conditions into the symmetry generator.
\item 
Integrate by parts in the dot product of the modified symmetry generator and the string equations to obtain the explicit expressions for the conserved density and its flux.
\end{enumerate}

\newpage 

\begin{center}
\textbf{Variational symmetries}
\end{center}
\begin{itemize}
\item
Case 1: Arbitrary $k(\lambda)$\\
In this case, there are 12 symmetries, of which 11 are variational symmetries. The uniform scaling is not variational.

\item
Case 2: $k(\lambda) = k_0 (1+\lambda)^{p}$,
\quad
$p \neq -4, -1, 0$, 
\quad
$k_0 = \const \neq 0$\\
In this case, there are 13 symmetries, of which 12 are variational. A combination of the scaling symmetries given by
\[\begin{aligned}
\vec{P} &= 
(
\frac{1}{p+4} ( -(3p + 4) t u_t + p (x+u)) - x (1 + u_x), \\
& \quad \quad \frac{1}{p+4} ( -(3p + 4) t v_t + p v) - x v_x, \\
& \quad \quad \frac{1}{p+4} ( -(3p + 4) t w_t + p w) - x w_x
)
\end{aligned}\]
is a variational symmetry.

\item
Case 3: $k(\lambda) = k_0 (1+\lambda)^{-4}$,
\quad
$k_0 = \const\neq 0$\\
In this case, there are 14 symmetries, of which 13 are variational. The uniform scaling symmetry is not variational.
\end{itemize}

\newpage

\begin{center}
\textbf{Conserved densities and their fluxes}
\end{center}

The variational symmetries in each case gives the following conservation laws. 
The physical meaning of the conserved densities $T_6, T_7, T_8, T_{12}, T_{13}, \text{ and } T_{14}$ is not known.

\begin{itemize}
\item
Case 1: Arbitrary $k(\lambda)$

In this case, there are 11 conservation laws corresponding to the 11 variational symmetries:
\begin{enumerate}[label=\roman*.]
\item 
time translation: $\vec{P} = (-u_t, -v_t, -w_t)$ gives 
\[\begin{aligned} 
T_1 &= \frac{1}{2} \rho A (u_t^2 + v_t^2 + w_t^2) + W(\lambda) = H, \\
X_1 &= - \frac{W'(\lambda)}{\lambda + 1} 
(
(1 + u_x) u_t + v_x v_t + w_x w_t
)
\end{aligned}\]
where $T_1$ is the \myemph{energy} density.
\item 
space translation: $\vec{P} = (-u_x, -v_x, -w_x)$ gives 
\[\begin{aligned}
T_2 &= \rho A (u_t u_x + v_t v_x + w_t w_x), \\
X_2 &= -\frac{1}{2} \rho A (u_t^2 + v_t^2 + w_t^2)
+ W(\lambda) \\
& \quad - \frac{W'(\lambda)}{\lambda + 1} 
(
(1 + u_x) u_x + v_x^2 + w_x^2
)
\end{aligned}\]
\item 
constant shift in $u$:
$\vec{P} =(1,0,0)$ 
gives 
\[
T_3 = u_t,
\quad 
X_3 = -k(\lambda) (1 + u_x)      
\]
where $T_3$ is the \myemph{$x$-momentum} density.
\item 
constant shift in $v$:
$\vec{P} =(0,1,0)$ 
gives
\[
T_4 = v_t,
\quad 
X_4 = -k(\lambda) v_x     
\]
where $T_4$ is the \myemph{$y$-momentum} density.
\item 
constant shift in $w$: 
$\vec{P} =(0,0,1)$ 
gives
\[
T_5 = w_t,
\quad 
X_5 = -k(\lambda) w_x     
\]
where $T_5$ is the \myemph{$z$-momentum} density.
\item 
time dependent shift in $u$: 
$\vec{P} =(t,0,0)$ 
gives 
\[
T_6 = u - t u_t,
\quad
X_6 =  t k(\lambda) (1 + u_x)
\]
\item 
time dependent shift in $v$: 
$\vec{P} =(0,t,0)$ 
gives 
\[
T_7 = v - t v_t,
\quad
X_7 = t k(\lambda) v_x
\]
\item 
time dependent shift in $w$: 
$\vec{P} =(0,0,t)$ 
gives 
\[
T_8 = w - t w_t,
\quad
X_8 = t k(\lambda) w_x
\]
\item 
rotations in the $(v,w)$ plane: $\vec{P} =(0,-w,v)$ gives 
\[
T_9 = v w_t - v_t w,
\quad
X_9 = k(\lambda) (v_x w - v w_x)
\]
where $T_9$ is the \myemph{angular momentum density in the $x$-direction}.
\item 
rotations in the $(x+u,w)$ plane: $\vec{P} =(w,0,-(x+u))$ gives 
\[
T_{10} = u_t w - (x+u) w_t,  
\quad
X_{10} = k(\lambda) ((x + u) w_x - (1 + u_x) w)
\]
where $T_{10}$ is the \myemph{angular momentum density in the $y$-direction}.
\item 
rotations in the $(x+u,v)$ plane: $\vec{P} =(-v,x+u,0)$ gives 
\[
T_{11} = (x+u) v_t - u_t v,
\quad
X_{11} = k(\lambda) ((1 + u_x) v - (x + u) v_x )
\]
where $T_{11}$ is the \myemph{angular momentum density in the $z$-direction}.
\end{enumerate}

\newpage
\item
Case 2: 
$k(\lambda) = k_0 (1+\lambda)^{p}$,
\quad
$p \neq -4, -1, 0$, 
\quad
$k_0 = \const \neq 0$
\\
In this case, corresponding to the 12 variational symmetries, there are 12 conservation laws, including specializations of the 11 from Case 1 plus one extra:
\begin{enumerate}[label=\roman*.]
\item
scaling: $\vec{P} = 
(
\frac{1}{p+4} ( -(3p + 4) t u_t + p (x+u)) - x (1 + u_x),
\frac{1}{p+4} ( -(3p + 4) t v_t + p v) - x v_x, 
\frac{1}{p+4} ( -(3p + 4) t w_t + p w) - x w_x
)$ 
gives 
\[\begin{aligned}
T_{12} &=
\frac{3p +4}{p + 4} t H
+ 
\rho A
(
x 
(u_t (1 + u_x) + v_t v_x + w_t w_x) 
\\
& \quad
- \frac{p}{p + 4}  
((x + u) u_t + v v_t + w w_t)
), \\
X_{12} &=
\rho A 
(
- \frac{1}{2} x (u_t^2 + v_t^2 + w_t^2)
- \frac{p + 1}{p + 2} k_0 x (1 + \lambda)^{p + 2} \\
& \quad + \frac{1}{p+4} k_0 (1 + \lambda)^p 
(
- (3p + 4) t(u_t (1 + u_x) + v_t v_x + w_t w_x) \\
& \quad + p ( (x + u) (1 + u_x) + v v_x + w w_x ) 
)
)
\end{aligned}\]
\end{enumerate}

\newpage
\item
Case 3: 
$k(\lambda) = k_0 (1+\lambda)^{-4}$,
\quad
$k_0 = \const\neq 0$


In this case, corresponding to the 13 variational symmetries, there are 13 conservation laws, including the specializations of the 11 from Case 1 plus two additional ones:
\begin{enumerate}[label=\roman*.]
\item 
scaling: $ \vec{P} = (\frac{1}{2} (x + u) - t u_t,
\frac{1}{2} v - t v_t,
\frac{1}{2} w - t w_t) $ gives
\[\begin{aligned}
T_{13} &= 
t H
-\frac{1}{2} \rho A ((x + u) u_t + v v_t + w w_t), \\
X_{13} &=  
\rho A k_0 (1 + \lambda)^{-4}
(
- t (u_t (1 + u_x) + v_t v_x + w_t w_x) \\ 
& \quad + \frac{1}{2} ((x + u) (1 + u_x) + v v_x + w w_x) 
)
\end{aligned}\]
\item 
time-dependent (conformal) scaling: $ \vec{P} = (t (x + u) - t^2 u_t,
t v - t^2 v_t,
t w - t^2 w_t)$ gives 
\[\begin{aligned}
T_{14} &= 
t^2 H
+ \rho A 
(
- t ((x + u) u_t + v v_t + w w_t)
+ \frac{1}{2} ((x + u)^2 + v^2 + w^2) 
), \\
X_{14} &= 
\rho A k_0 t (1 + \lambda)^{-4} 
(
- t  (u_t (1 + u_x) + v_t v_x + w_t w_x) \\
& \quad + (x + u) (1 + u_x) + v v_x + w w_x
)
\end{aligned}\]
\end{enumerate}
\end{itemize}

\newpage
\begin{center}
\textbf{Conclusions}
\end{center}

\begin{itemize}
\item
\myemph{Derivations of the equations of motion for an elastic string}
\\
The string equations can be derived from three different perspectives: 
Newtonian mechanics, nonlinear elasticity, and continuum mechanics. 
Each approach is consistent with the others.
\item 
\myemph{Symmetries}
\\
We found a 12-dimensional symmetry group for the string equations when the stress-strain function is arbitrary.
There is one extra scaling symmetry in the power case and two additional symmetries: 
the scaling and a time-dependent scaling, in the special power case.
\item
\myemph{Conservation laws}
\\
Using Noether's theorem, we found the conservation laws admitted by the string equations. They include the familiar conserved quantities: energy, um, and angular momenta.
\newpage
\item
\myemph{No travelling wave solutions}
\\
A direct substitution of the travelling wave solution ansatz $\vec{r} = \vec{r}(x - ct)$ into the string equations, followed by integrating once, shows that travelling wave solutions do not exist. 
\item 
\myemph{Non-straight equilibrium configurations}
\\
The equilibrium condition for an arbitrary string configuration with arc length parameter $s$, $\partial_s (T(s)\widehat{r_s}) = 0$, gives two solutions:
\\
1. a straight line configuration with constant tension, or
\\
2. an arbitrary curve configuration with zero tension. 
\item  
\myemph{Prospective areas of research:}
\\
1. finding the symmetries and conservation laws in the case of a finite string with physical boundary conditions at the ends 
\\
2. generalizing the derivation of the string equations to membranes, the two-dimensional analog of strings.
\end{itemize}
















% References

% W.F. Osgood, 
% Advanced Calculus, 
% The Macmillian Company, New York, 
% 1925. 

% G.F. Carrier, 
% On the nonlinear vibration problem of the elastic string. 
% Q. J. Appl. Math. 3 (1945), 157--165.

% R. Narasimha, 
% Nonlinear vibration of an elastic string, 
% J. Sound and Vibration 8 (1968), 134--146.

% O. Serrano, R. Zaera, J. Fern\'andez-S\'aez, M. Ruzzene, 
% Generalized continuum model for the analysis of nonlinear vibrations of taut strings with microstructure, 
% Internat. J. Solids and Structures 164 (2019),  157--167. 

% S.S. Antman, 
% Noninear Problems of Elasticity, 
% (Applied Math. Sci., volume 107), 
% Springer, 2005. 

% A.S. Vasudeva Murthy,
% On the string equation of Narasimha, 
% In: Connected at infinity II: a selection of mathematics by Indians, 
% Eds. R. Bhatia, C.S. Rajan, A.I. Singh
% (TRM, volume 67), 58--84. 
% Hindustan Book Agency, 2013.


\end{document}
